\documentclass[11pt]{article}
\usepackage[margin=1in]{geometry}
\usepackage{amsmath,amssymb,booktabs,hyperref}
\usepackage[T1]{fontenc}
\newcommand{\Oh}{\mathrm{Oh}}
\newcommand{\We}{\mathrm{We}}
\newcommand{\Bo}{\mathrm{Bo}}
\title{The contact-edge selector: why $\theta_c=\inf\mathcal F$ fails,\\
{\large and the closed-form tangency residual that replaces it}}
\author{SpectralKM derivations}
\date{}

\begin{document}
\maketitle

\section{The symptom --- retracted}

\begin{quote}
\textbf{An earlier version of this note opened by asserting that the model overpredicts the
contact time by roughly a factor of two. That is retracted: it was an artefact of the comparison,
not a property of the model.}
\end{quote}

Run at \emph{identical} dimensionless parameters, this model and the reference implementation
(harrislab-brown/BouncingDroplets, executed directly in \textsc{Matlab}) agree:

\begin{center}
\begin{tabular}{lccl}
\toprule
metric & ours & reference & difference \\
\midrule
contact duration (\path{primary_contact_time} / their $f>0$) & $3.6293$ & $3.6530$ & $0.65\%$ \\
threshold, centre of mass below $R$ & $4.4432$ & $4.5297$ & $1.9\%$ \\
$\min z_{cm}/R$ & $0.1670$ & $0.1110$ & we penetrate less \\
\bottomrule
\end{tabular}
\end{center}

at $\We=1.0959$, $\Bo=0.01661$, $\Oh=0.00617$ ($R=0.35$\,mm water, $U=0.4759$\,m/s).

\paragraph{Where the false symptom came from: the wrong physical problem.} The dataset the
comparison was built on is the \textbf{droplet-onto-solid-substrate} problem, not the
droplet-onto-bath problem this model solves. The spreadsheet is named after the repository
\path{LowWeberDropRebound}, whose \path{README} opens ``MATLAB scripts for simulating the dynamics
of liquid droplets impacting \emph{solid substrates}''. Six independent checks agree:

\begin{center}
\begin{tabular}{ll}
\toprule
Check & Result \\
\midrule
that repo's \path{sweeper_experiments.m} & hardcodes \verb|Bo = 0.0189| --- exactly the single \\
 & fixed $\Bo$ of the file's \texttt{dns} and \texttt{km\_model} blocks \\
\verb|besselj|, \verb|tanh|, \verb|bath| in its solver & \textbf{zero} occurrences: no bath surface
degrees \\
 & of freedom at all, only the drop's $\beta_l$ \\
the bath paper's Figure 5 & water at $\Bo=0.017$, $\Oh=0.006$, with
quasi-potential \\
 & and DNS curves \emph{only} --- no ``KM model'' curve \\
the bath paper's experiments & exactly \emph{two} $(\Bo,\Oh)$ points (water, oil); \\
 & this file's is a continuum, $\Oh=0.0139\ldots0.7865$ \\
the file's \texttt{km\_model} block & includes $\Bo=0$ and $\Oh=0$ rows, i.e.\ an inviscid \\
 & zero-gravity sweep, not a two-fluid experiment \\
contact-time bands & disjoint (below) \\
\bottomrule
\end{tabular}
\end{center}

The last check is the physical one, and it is decisive. On a \emph{rigid wall} the rebound is set
by the drop's own $l=2$ free oscillation, full period $2\pi/\sqrt{8}=2.2214\,t_\sigma$, with the
classical superhydrophobic value $\approx2.6\,t_\sigma$ \cite{RichardClanetQuere2002}; that file
gives $\texttt{km\_model}=2.41$ at $\We=1$ and $\texttt{dns}\in[1.96,3.36]$. A \emph{bath} deforms
and carries impact energy away as interfacial waves, so the drop rides down and back up and
contact lasts markedly longer: the bath dataset (\path{lowWeberComparison.csv},
$153$ runs at $N=60$, $89$ with a measured $t_c$) spans $[4.4408,8.0946]$. The two bands do not
overlap, so no amount of parameter matching reconciles them --- they are answers to different
questions. This is now pinned in \path{test/test_bath_reference.jl}.

At the bath reference's own water point, $\Bo=0.016644$, $\Oh=0.006165$, it reports
$t_c=4.5286$; running the reference \textsc{Matlab} directly gives $4.5297$, and this model gives
$4.4432$ --- squarely in the bath band. The earlier ``puzzle'' of $4.53$ against a
\texttt{km\_model} value of $2.41$, which appeared to contradict the paper's own claim that $t_c$
is nearly $\Oh$-independent, was never a puzzle: those are a bath and a rigid wall.

\paragraph{Against experiment, not against another code.} All of the above is
model-against-model, which settles the problem identity and nothing else. The bath
\emph{experiments} are now in the repository as \path{bath_experiment_water.csv} and
\path{bath_experiment_oil.csv}: $17$ measured points with standard deviations over at least five
trials each, extracted by \path{scripts/extract_bath_experiment.m} from the authors' own
\path{.fig} files. They were absent from the digitised CSV export that already existed, because
the experimental points are stored as \texttt{errorbar} objects --- one per point --- and an
export that walks only \texttt{line} children drops every one of them; that export holds the
figure's eight model curves and none of its five water points.

\begin{center}
\begin{tabular}{lccc}
\toprule
 & $\We$ & $t_c/t_\sigma$ & \\
\midrule
experiment (water, $\Bo=0.017$, $\Oh=0.006$) & $0.725$ & $4.6628\pm0.1887$ & $\ge5$ trials \\
this model, \path{threshold_contact_time}    & $0.700$ & $4.5039$ & \\
\midrule
residual & & $-3.4\%$ & $0.84\,\sigma$ \\
\bottomrule
\end{tabular}
\end{center}

Two things keep that number honest. The metrics are \emph{not} the same: the experiment times the
north pole across $z=2R$, detachment not being optically resolvable, whereas the model convention
--- which \path{threshold_contact_time} implements --- is the centre of mass across $z=R$, and the
authors quote a typical $5\%$ difference between the two for water. And the paper states that its
own quasi-potential model and DNS ``slightly underpredict the dimensionless contact time at
intermediate $\We$'', attributing it to neglected nonlinear deformation. A residual of $-3.4\%$ is
therefore consistent with the definition offset alone, and its sign is the one the authors report.
Pinned in \path{test/test_bath_reference.jl}.

\paragraph{A second, compounding error.} Independently of the wrong-problem error, the comparison
script selected parameters by \emph{binning on $\We$ alone}: it split $\log_{10}\We$ into equal
intervals over the $931$ points and took the median $\We$, $\Bo$ and $\Oh$ of whatever fell in
each, so $\Bo$ and $\Oh$ were never matched to any actual measurement. A single window mixes
fluids differing by $33$--$54\times$ in $\Oh$, and the $t_c$ values \emph{within} one window
already differ by $1.35$--$1.48\times$. That confound was noted in writing when the script was
created and not acted upon for several rounds. It is fixed --- \path{choose_points} now clusters
on $(\We,\Bo,\Oh)$ jointly and reports each point's $\Oh$ and $\Bo$ spread --- but it was the
smaller of the two mistakes.

\paragraph{What was eliminated along the way.} Each by measurement, and each still worth knowing:

\begin{center}
\begin{tabular}{lll}
\toprule
Candidate & Test & Outcome \\
\midrule
metric definition & \path{threshold_contact_time} vs.\ \path{contact_time} & real, and it does differ \\
drop viscous closure & \texttt{:reid} vs.\ \texttt{:lamb} & 4th digit \\
bath radius (reflected waves) & \path{debug_bath_radius_sensitivity.jl} & flat: $4.9402$, $4.9613$, \\
 & $b=6,12,20,30$ with $M\propto b$ & $4.9614$, $4.9614$ \\
contact-edge rule & \texttt{:crossing} (the reference's own rule) & unchanged to 3 digits \\
\bottomrule
\end{tabular}
\end{center}

$b=6$ remains unvalidated per \path{provenance.tex}; it is now measured as \emph{not} affecting
the contact time.

\section{What is actually wrong}

Tracing a contact (\path{debug_contact_exit_trace.jl}, \path{debug_feasibility_collapse.jl}):

\begin{itemize}
  \item the net force $f\le0$ on \textbf{0 of 221} steps -- it decays asymptotically to
        $1.4\times10^{-8}$ from above and never reverses, so the $f\le0$ exit test never fires
        and contact terminates only at the $\theta_c$ floor;
  \item $63\%$ of the contact duration occurs \emph{after} the droplet has begun rising;
  \item at $t=0.009$ the non-intersection predicate is false for all $\theta<0.0492$ and
        $\inf\mathcal F=0.0492$; one step later, at $t=0.010$, it is true even at
        $\theta=10^{-4}$, $\inf\mathcal F$ collapses by a factor of $200$, and the patch never
        recovers.
\end{itemize}

\emph{None of this explains a contact-time error, because there is none to explain.} What follows
stands on its own: the rule is degenerate, that degeneracy is real and reproducible, and it
produces a physically indefensible patch at low $\We$. It is a defect worth fixing for its own
sake, not the cause of a symptom that turned out to be an artefact.

The search is not at fault. A brute-force scan (\path{debug_selector_rate_limit.jl}) shows
\path{select_theta_c} returning within $0.3$--$4\%$ of the true $\inf\mathcal F$ at every step.
The \emph{rule} is at fault.

\section{The degeneracy, stated precisely}
\label{sec:degeneracy}

Write the gap as $C(\theta,\tau)=\eta(r(\theta,\tau),\tau)-z_{cm}+z_d(\theta,\tau)$. The
kinematic match imposes $C=0$ on $[0,\theta_c]$; physical consistency requires $C<0$ beyond the
edge. Continuing just outside, $C(\theta)\simeq a\,(\theta-\theta_c)$ with
\begin{equation}
  a(\theta_c) := \partial_\theta C\big|_{\theta_c}\text{,}
\end{equation}
so non-overlap is the inequality $a(\theta_c)\le0$ and the feasible set is
$\mathcal F=\{\theta_c: a(\theta_c)\le0\}$. The two rules are then
\begin{equation}
  \text{current: } \theta_c=\inf\mathcal F\text{,}\qquad
  \text{tangency \eqref{eq:tangency}: } a(\theta_c)=0\text{.}
\end{equation}

\begin{quote}
If $a(\theta_c)<0$ for every $\theta_c>0$ -- the measured situation from $t=0.010$ onward --
then $\mathcal F=(0,\theta_{\max}]$ and $\inf\mathcal F=0$, \emph{regardless of the shape or
magnitude of $a$}.
\end{quote}

This is the whole failure. The infimum rule uses only the \textbf{sign} of $a$ and never its
magnitude, so it returns the same answer whether the surfaces are barely separated or widely so.
It carries information only while the constraint is \emph{binding}; the moment non-intersection
goes inactive, it degenerates to $\theta_c=0$ -- not because the physics says the patch vanishes,
but because the rule has nothing left to determine it with. Verified symbolically in
\path{derivations/cas_tangency_residual.py}, which also exhibits an $a<0$ example where
$\inf\mathcal F=0$ while $\arg\min|a|$ is interior.

Two further facts from the scan bear on the rule's stated justification. Positivity ($p>0$) holds
\emph{at} the collapsed patch, so complementarity -- which \path{select_theta_c}'s docstring cites
as its basis -- does not forbid it. And the band's upper edge is set by $f<0$ and $p<0$, i.e.\ by
suction, not by \texttt{eq:check-monotone-r}. The rule's original evidence was a single measurement
at one step at onset, precisely where the constraint \emph{is} active and the rule does work.

\section{The closed-form tangency residual}

\texttt{eq:tangency-selector} is $T=\partial_\theta C|_{\theta_c}$, currently evaluated in
\path{src/residual.jl} by a centred finite difference with $h=10^{-6}$.

\paragraph{A claim withdrawn.} An earlier draft of this note asserted that differencing a
Fourier--Bessel sum at $h=10^{-6}$ ``discards roughly half the available digits'' and was
therefore unfit to root-find. \emph{That was wrong, and the test written to check it said so}:
against a Richardson-extrapolated reference at $h=10^{-3}$, the existing finite-difference
residual agrees to $\texttt{rtol}=10^{-8}$ at $\theta=0.25,0.6,1.1$ (\path{test/test_selector.jl}).
Accuracy is not the problem. The closed form below is still worth having --- it removes a
step-size parameter, and it differentiates cleanly under \path{ForwardDiff} where a difference of
a difference would not --- but the case for it is structural, not numerical, and the fix to the
selector does not depend on it. With
$\xi(\theta)=1+\sum_l\beta_lP_l(\cos\theta)$, $r=\xi\sin\theta$, $z_d=-\xi\cos\theta$ and
$\eta(r)=\sum_ma_mJ_0(k_mr)$, the chain rule gives
\begin{equation}
  \boxed{\;T=\eta'(r)\,\big[\xi'\sin\theta+\xi\cos\theta\big]
           +\big[\xi\sin\theta-\xi'\cos\theta\big]\;}
  \label{eq:tangency}
\end{equation}
with
\begin{equation}
  \xi'(\theta)=-\sin\theta\sum_l\beta_l P_l'(\cos\theta)\text{,}\qquad
  \eta'(r)=-\sum_m a_m k_m J_1(k_m r)\text{.}
\end{equation}
\eqref{eq:tangency} is verified against symbolic differentiation of the composed expression
(difference exactly zero), and against the existing numerical residual in
\path{test/test_selector.jl}.

\section{The pole degeneracy is physical, not algebraic}

$T\equiv0$ at $\theta=0$ and $\theta=\pi$ for any state -- \texttt{eq:tangency-degenerate}. Worth
recording \emph{how} the CAS established this, because the first attempt was wrong: with $\eta$
left as an abstract function the CAS correctly \textbf{refused}, returning
$T(0)=\xi(0)\,\eta'(0)$, since nothing in the chain rule makes $\eta'(0)$ vanish. It vanishes
because of \textbf{axisymmetry}, through $J_1(0)=0$; and $\xi'$ carries a factor $\sin\theta$, so
it vanishes at both poles too. Substituting the concrete expressions, $T\to0$ at both poles for
$l=2$ and $l=3$ and arbitrary $\beta_l,a_m,k_m,z_{cm}$.

Two implementation consequences follow. A seed is required at onset, which
\path{onset_theta_c} already supplies from the geometric crossing $C=0$; and the residual must be
root-found \emph{away} from $\theta=0$, never bracketed from it.

\section{The replacement rule}

Select $\theta_c$ as a root of \eqref{eq:tangency}, retaining
\texttt{eq:check-nonintersect} and \texttt{eq:check-monotone-r} as a \emph{feasibility filter} over
candidates -- which is what \path{paper-formulation.tex} describes, and what
\cite{AlventosaEtAl2023} and \cite{GaleanoRios2017} do. Because $\eta(r(\theta))$ is a sum of
oscillatory $J_0$ terms, $T$ can change sign more than once on $(0,\pi)$, so a plain bisection on
a sign change is not safe; the argmin form of \texttt{eq:theta-c-argmin} over admissible candidates
is used instead, seeded from the previous step and, at onset, from \path{onset_theta_c}.

\section{The degeneracy is fatal at high Weber number: \texttt{:crossing} becomes the default}
\label{sec:crossing-default}

An earlier revision of this note closed by saying that the degeneracy, though real, changed no
observable, and \path{DEFAULT_SELECTOR} was therefore left at \texttt{:feasible}. \textbf{That
verdict was drawn from too narrow a sample and is retracted here.} It rested on two low-Weber water
points, $\We=0.0231$ and $\We=0.9985$, where the two rules agree to three digits. They do. At high
$\We$ they do not agree at all: one returns a physical trajectory and the other returns a droplet
that sinks.

At $\We=7.307$ in $5$~cSt silicone oil ($\Bo=0.0563$, $\Oh=0.0578$), \texttt{:feasible} holds
contact for $21$ steps --- a duration of $0.194$ against a measured $5.26$ --- then detaches, after
which the droplet free-falls to $z_{cm}=-84.2$ by $\tau=30$ and never re-contacts. No contact time
exists to report. Under \texttt{:crossing} the same case gives contact over $(0,4.229)$, a rebound
to $z_{cm}=2.239$, and further contacts at $\tau\in(18.33,22.15)$ and $\tau>29.38$: a physical
multi-bounce trajectory. This is the degeneracy of \S\ref{sec:degeneracy} at full strength --- once
non-intersection stops binding, $\inf\{\theta:\ \text{feasible}\}$ is $0$ regardless of the gap, and
at high $\We$ that happens almost immediately rather than late in contact.

\subsection{The change is a strict improvement, measured}

The switch was gated on both observables the experiments report, at the default truncation
($M=60$, $L=120$, $N=3$, $n_q=200$), at five anchor points spanning both working fluids, with each
rule run on identical hardware. Contact time is $t_c/t_\sigma$; $\delta$ is the maximum penetration
depth of the droplet's south pole, taken over the \emph{first} contact interval.

\begin{center}
\begin{tabular}{llrrrr}
\hline
fluid & $\We$ & $t_c$ \texttt{:feasible} & $t_c$ \texttt{:crossing} & $\delta$ \texttt{:feasible} & $\delta$ \texttt{:crossing} \\
\hline
water & 0.7251 & 4.49560 & 4.49657 & 0.66951 & 0.66900 \\
water & 1.9387 & 4.42119 & 4.42242 & 0.95812 & 0.95736 \\
oil   & 1.2158 & 5.07868 & 5.08034 & 0.72344 & 0.72301 \\
oil   & 3.9463 & 5.01559 & 5.01686 & 1.07481 & 1.07417 \\
oil   & 7.3070 & \emph{none} & 5.03505 & 0.30009 & 1.33536 \\
\hline
\end{tabular}
\end{center}

Wherever \texttt{:feasible} produces an answer at all the two rules agree to four significant
figures --- at most $0.03\%$ in $t_c$ and $0.08\%$ in $\delta$ --- so \texttt{:crossing} cannot be
said to regress any case that previously worked. Where \texttt{:feasible} fails it fails completely.
That is a strict improvement, and it is why the default changes. The cost of the change to stored
results is correspondingly small: any low-$\We$ number regenerated after this commit moves in its
fourth digit, and any high-$\We$ oil number changes qualitatively because it was not physical
before.

\subsection{What the gate did \emph{not} vindicate: $\delta$}

The same measurement exposes a separate deficit that is \emph{not} selector-related, since it is
identical to four figures under both rules. Against the published quasi-potential model (1PKM),
interpolated from the same figures the experimental points were read from:

\begin{center}
\begin{tabular}{llrrr}
\hline
fluid & $\We$ & $\delta$ here & $\delta$ 1PKM & $\delta$ measured \\
\hline
oil & 1.2158 & 0.7230 ($-1.60\sigma$) & 0.7415 ($-1.30\sigma$) & $0.8229\pm0.0626$ \\
oil & 3.9463 & 1.0742 ($-4.69\sigma$) & 1.2811 ($-2.65\sigma$) & $1.5495\pm0.1013$ \\
oil & 7.3070 & 1.3354 ($-4.40\sigma$) & 1.7042 ($-2.15\sigma$) & $2.0558\pm0.1637$ \\
\hline
\end{tabular}
\end{center}

Under-prediction of $\delta$ is a property of the reduced-model class, not of this package: over the
twelve oil points 1PKM averages $-1.64\sigma$ ($-12.3\%$) and clears $2\sigma$ at only $8$ of $12$,
while the DNS averages $-0.71\sigma$ and clears it at all $12$. But this package under-predicts
$\delta$ roughly twice as far as 1PKM does at moderate and high $\We$, while simultaneously
predicting $t_c$ \emph{better} than 1PKM there ($-0.85\sigma$ against $-3.61\sigma$ at
$\We=7.307$). A model that gets the contact duration right and the penetration depth wrong is
saying something specific about where its linearisation fails, and that is open.

\section{Scope}

This note establishes that the current rule is degenerate, and \S\ref{sec:crossing-default}
establishes that the degeneracy destroys the high-$\We$ solution outright, which is why
\texttt{:crossing} is now the default. What remains open is the $\delta$ deficit above, which is
selector-independent and roughly twice 1PKM's at moderate $\We$. Any future claim here must be made
against data for the \emph{same physical problem} --- a bath, not a solid substrate --- and then
\emph{parameter-matched}, on $\Bo$ and $\Oh$ and not $\We$ alone. Nor does this note
revisit the upper edge of the admissible band, where $f<0$ and $p<0$ already exclude wide patches;
whether that exclusion is itself too aggressive is a separate question, and the $\delta$ deficit
makes it a more interesting one --- too aggressive an upper edge would under-predict penetration.

\begin{thebibliography}{9}
\bibitem{AlventosaEtAl2023} L.~F.~L. Alventosa, R. Cimpeanu, D.~M. Harris,
  \emph{Inertio-capillary rebound of a droplet impacting a fluid bath}, JFM (2023).
\bibitem{GaleanoRios2017} C.~A. Galeano-R\'ios, P.~A. Milewski, J.-M. Vanden-Broeck,
  \emph{Non-wetting impact of a sphere onto a bath}, JFM (2017).
\bibitem{RichardClanetQuere2002} D. Richard, C. Clanet, D. Qu\'er\'e,
  \emph{Contact time of a bouncing drop}, Nature \textbf{417}, 811 (2002).
\end{thebibliography}

\end{document}
